\section{WSL}

\subsection{安装WSL}

管理员模式下打开 PowerShell 或 Windows 命令提示符，然后(需要翻墙)运行：

\begin{lstlisting}
  wsl --install
\end{lstlisting}

\subsection{安装Texlive}

\begin{enumerate}
  \item 下载安装镜像
  
  清华源 \link{https://mirrors.tuna.tsinghua.edu.cn/CTAN/systems/texlive/Images/}
  \item 挂载安装镜像
  
  \begin{lstlisting}
    sudo mkdir /mnt/texlive
    sudo mount /mnt/d/texlive2023.iso /mnt/texlive    
  \end{lstlisting}
  \item 启动安装程序
  
  \begin{lstlisting}
    sudo /mnt/texlive/install-tl % 输入I并回车，开始执行安装。
  \end{lstlisting}
  \item 安装善后
  
  \begin{lstlisting}
    sudo umount /mnt/texlive	% 将安装镜像弹出，注意不是unmount! 
    sudo rm -r /mnt/texlive % 删除/mnt/texlive文件夹
  \end{lstlisting}
  \item 配置环境变量
  
  \begin{lstlisting}
    sudo apt-get update % 更新apt
    sudo apt install gedit % 安装gedit
    sudo gedit ~/.bashrc
    % 在.bashrc文件的末尾添加如下代码
    export PATH=/usr/local/texlive/2023/bin/x86_64-linux:$PATH
    export MANPATH=/usr/local/texlive/2023/texmf-dist/doc/man:$PATH
    export INFOPATH=/usr/local/texlive/2023/texmf-dist/doc/info:$PATH
    source ~/.bashrc % 重新加载 .bashrc 文件
    $
  \end{lstlisting}
  \item 重启设备
  \item 查看是否安装成功
  
  \begin{lstlisting}
    tex --version
  \end{lstlisting}
\end{enumerate}

\subsection{安装windows字体}

\begin{enumerate}
  \item 将windows下进入C盘的 C:/windows/Fonts拷贝到D:/winfonts
  \item 新建目录用于存储字体
  
  \begin{lstlisting}
    sudo mkdir /usr/share/fonts/winfonts
  \end{lstlisting}
  \item 将windows系统下的字体拷贝到刚刚创建的winfonts目录下
  
  \begin{lstlisting}
    sudo cp -r /mnt/d/winfonts/* /usr/share/fonts/winfonts
  \end{lstlisting}
  \item 修改字体文件的权限，使root用户以外的用户也可以使用
  
  \begin{lstlisting}
    cd /usr/share/fonts/
    sudo chmod -R 755 winfonts
  \end{lstlisting}
  \item 将texlive配置文件复制到系统
  \begin{lstlisting}
    sudo cp /usr/local/texlive/2023/texmf-var/fonts/conf/texlive-fontconfig.conf /etc/fonts/conf.d/09-texlive.conf
  \end{lstlisting}
  \item 安装mkfontscale mkfontdir和fc-cache命令
  
  \begin{lstlisting}
    % 使mkfontscale和mkfontdir命令正常运行
    sudo apt-get install ttf-mscorefonts-installer
    % 使fc-cache命令正常运行
    sudo apt-get install fontconfig
  \end{lstlisting}
  \item 建立字体缓存
  
  \begin{lstlisting}
    mkfontscale
    mkfontdir
    fc-cache -fv
  \end{lstlisting}
\end{enumerate}

\subsection{VSCode配置}

\begin{itemize}
  \item hsnips文件
  
  \begin{lstlisting}
    sudo mkdir -p ~/.config/Code/User/hsnips
    sudo cp -r /mnt/c/Users/bz/AppData/Roaming/Code/User/hsnips/* ~/.config/Code/User/hsnips
  \end{lstlisting}
  \item 模版
  
  \begin{lstlisting}
    sudo cp /usr/local/texlive/2023/texmf-var/fonts/conf/texlive-fontconfig.conf /etc/fonts/conf.d/09-texlive.conf
    sudo mkdir -p ~/texmf/tex/latex
    sudo cp -r /mnt/c/Users/bz/texmf/tex/latex* /home/bz/texmf/tex
  \end{lstlisting}
\end{itemize}

\subsection{Gitee同步}

\begin{enumerate}
  \item 设置帐号
  
  \begin{lstlisting}
    git config --global user.name "zhangbin"
  \end{lstlisting}
  \item 设置邮箱
  
  \begin{lstlisting}
    git config --global user.email "12267558+vmz@user.noreply.gitee.com"
  \end{lstlisting}
  \item 生成密钥
  
  \begin{lstlisting}
    ssh-keygen -t rsa -C "12267558+vmz@user.noreply.gitee.com"
    % 三次回车即可生成 ssh key
  \end{lstlisting}
  \item 查看生成的公钥
  
\begin{lstlisting}
    cat ~/.ssh/id_rsa.pub
\end{lstlisting}
  \item 复制至gitee
  
  打开gitee官网，设置->SSH公钥->添加 SSH公钥。

  标题可以任意设置，生成的公钥\verb|id_rsa.pub|内容复制到gitee，就可以添加SSH公钥了。
  \item 测试验证是否连接成功
  
  \begin{lstlisting}
    ssh -T git@gitee.com
    % 如出现以下欢迎信息, 说明连接成功, 进一步证明上面的配置正确. 否则, 请仔细检查上面的步骤.
    Hi xxx! You've successfully authenticated, but GITEE.COM does not provide shell access.
  \end{lstlisting}
  \item 创建本地仓库
  
  \begin{lstlisting}
    sudo mkdir ./workspace
    cd workspace
  \end{lstlisting}
  \item 赋予目录账户权限
  
  \begin{lstlisting}
    sudo chown -R bz /home/bz/workspace/
  \end{lstlisting}
  \item 把这个目录变成Git可以管理的仓库
  
  \begin{lstlisting}
    git init
  \end{lstlisting}
  \item 把本地仓库与一个云端Gitee仓库关联（如果从已有云端仓库拉取直接进行下一步）
  
  \begin{lstlisting}
    git remote add origin https://gitee.com/vmz/latex.git
  \end{lstlisting}
  \item 拉取
  
  \begin{lstlisting}
    git pull https://gitee.com/vmz/latex.git master
  \end{lstlisting}
  \item VSCode每次同步都要输入用户名密码？
  
  在 VSCode 的终端输入 \verb|git config --global credential.helper store| 命令

  在弹框中输入账号和密码，此时输入一次，以后再git push /pull 的时候就不用再输账号和密码了。
\end{enumerate}

